\textbf{Многомерные обратные задачи для пар $(p, h)$ и $(p, s)$ в регионе~3.}
Если состояние в регионе~3 задаётся входной парой $(p, h)$
или $(p, s)$,
задача уже не сводится к одномерному поиску.
Неизвестными становятся сразу две переменные:
\begin{equation}
\mathbf{z} =
\left[
\begin{array}{c}
\rho \\
T
\end{array}
\right],
\qquad
\Phi(\mathbf{z}) = \frac{1}{2}\|\mathbf{r}(\mathbf{z})\|_{2}^{2}.
\end{equation}
Для пары $(p, h)$ вектор невязки имеет вид
\begin{equation}
\mathbf{r}(\mathbf{z}) =
\left[
\begin{array}{c}
p_{\mathrm{if97}}(\rho, T) - p_{\mathrm{target}} \\
h_{\mathrm{if97}}(\rho, T) - h_{\mathrm{target}}
\end{array}
\right],
\end{equation}
а для пары $(p, s)$
во второй компоненте используется
$s_{\mathrm{if97}}(\rho, T) - s_{\mathrm{target}}$.

Метод Ньютона--Рафсона
линеаризует систему в окрестности текущего приближения
и на каждом шаге решает задачу
\begin{equation}
\begin{aligned}
J_{k}\,\Delta \mathbf{z}_{k} &= -\mathbf{r}_{k}, \\
\mathbf{z}_{k+1} &= \mathbf{z}_{k} + \Delta \mathbf{z}_{k},
\end{aligned}
\end{equation}
\begin{gostwhere}
\item $J_{k} = \partial \mathbf{r}/\partial \mathbf{z}$ --- матрица Якоби.
\end{gostwhere}

При хорошем начальном приближении
метод обеспечивает быструю локальную сходимость.
Однако в околокритической области
он оказался недостаточно устойчивым:
из-за роста сжимаемости,
уплощения изотерм и вырождения производных давления по плотности
матрица $J_{k}$ становится плохо обусловленной,
а полный шаг $\Delta \mathbf{z}_{k}$
может выбрасывать траекторию за границы допустимого региона
или переводить её на неверную фазовую ветвь.

Поэтому применяется
схема Левенберга--Марквардта:
внутренний цикл решателя строится по тому же принципу,
но работает уже с вектором невязки, матрицей Якоби
и адаптивным демпфированием шага.
\begin{equation}
\left(J_{k}^{\mathsf{T}}J_{k} + \lambda_{k} I\right)\Delta \mathbf{z}_{k}
= -J_{k}^{\mathsf{T}}\mathbf{r}_{k},
\end{equation}
\begin{equation}
\mathbf{z}_{k+1} = \mathbf{z}_{k} + \alpha_{k}\Delta \mathbf{z}_{k},
\qquad
0 < \alpha_{k} \le 1.
\end{equation}
Здесь параметр $\lambda_{k}$ выполняет роль регуляризации:
при больших значениях метод ведёт себя ближе к градиентному спуску,
а при малых --- приближается к гаусс-ньютоновскому шагу.
Демпфирующий множитель $\alpha_{k}$
ограничивает длину шага
и предотвращает перескок через фазовые границы.

В реализованной модификации используются три дополнительных приёма.
Во-первых, начальное приближение выбирается
не произвольно,
а по результатам предварительного сеточного поиска:
\begin{equation}
\mathbf{z}_{0} = \arg\min_{\mathbf{z}\in\mathcal{G}}
\|\mathbf{r}(\mathbf{z})\|_{2},
\end{equation}
\begin{gostwhere}
\item $\mathcal{G}$ --- грубая сетка в допустимой области региона~3.
\end{gostwhere}

Во-вторых, параметр $\lambda_{k}$ адаптируется по качеству шага:
если норма невязки уменьшается,
то $\lambda_{k}$ снижается,
а если шаг оказывается неудачным ---
увеличивается.
В-третьих, для подбора $\alpha_{k}$
используется демпфирование шага,
например по правилу обратного деления:
\begin{equation}
\begin{aligned}
\alpha_{k} &= 2^{-m_{k}}, \\
m_{k} &= \min\left\{
m \ge 0 :
\|\mathbf{r}(\mathbf{z}_{k} + 2^{-m}\Delta \mathbf{z}_{k})\|_{2}
\right. \\
&\qquad \left.
<
\|\mathbf{r}(\mathbf{z}_{k})\|_{2}
\right\}.
\end{aligned}
\end{equation}
Совместно эти модификации решают три разные проблемы:
сеточный поиск уменьшает зависимость от стартовой точки,
регуляризация стабилизирует почти вырожденный Якобиан,
а демпфирование предотвращает осцилляции и выход за пределы региона~3.

На рисунке~\ref{fig:convergence_comparison} показана характерная динамика
сходимости невязки для двух типовых режимов работы решателя.
В однофазной области, удалённой от линии насыщения,
метод Ньютона--Рафсона достигает целевой точности
за несколько итераций,
тогда как в околокритической области
тот же контур переходит в демпфированный режим
и требует заметно большего числа шагов.
Именно это поведение определяет практическую структуру
итеративных решателей в проекте:
в стабильных областях используется быстрый локальный шаг,
а в сложных зонах включаются line search,
регуляризация и контроль региона.

\begin{figure}[htbp]
    \centering
    \begin{tikzpicture}
    \begin{axis}[
        width=\textwidth,
        height=8cm,
        ymode=log,
        ymin=1e-9, ymax=1e1,
        xmin=0, xmax=15,
        xlabel={Номер итерации, $N$},
        ylabel={Абсолютная невязка, $|f(x)|$},
        grid=major,
        grid style={dashed, gray!50},
        legend pos=north east,
        legend style={cells={anchor=west}, font=\footnotesize},
        mark options={solid}
    ]
        \addplot [domain=0:15, dashed, thick, red] {1e-7};
        \addlegendentry{Целевая точность ($10^{-7}$)}

        \addplot[
            color=blue,
            mark=square*,
            thick
        ] coordinates {
            (1, 2.5e0)
            (2, 4.1e-1)
            (3, 1.2e-2)
            (4, 3.5e-5)
            (5, 1.8e-9)
        };
        \addlegendentry{Регион 5 (однофазная область, 1D)}

        \addplot[
            color=green!60!black,
            mark=*,
            thick
        ] coordinates {
            (1, 3.0e0)
            (2, 1.5e0)
            (3, 8.0e-1)
            (4, 4.5e-1)
            (5, 2.0e-1)
            (6, 8.5e-2)
            (7, 3.1e-2)
            (8, 9.0e-3)
            (9, 1.5e-3)
            (10, 2.2e-4)
            (11, 1.1e-5)
            (12, 4.0e-7)
            (13, 1.2e-8)
        };
        \addlegendentry{Регион 3 (околокритическая область, 2D)}

    \end{axis}
    \end{tikzpicture}
    \caption{Сравнительная динамика сходимости итерационных решателей в однофазной и околокритической областях IF97}
    \label{fig:convergence_comparison}
\end{figure}
